\chapter{Đào sâu: File ở quy mô lớn --- presigned URL, endpoint, và vì sao byte không thuộc về Postgres}
\label{ch:dd-storage}

Chương~9 giới thiệu MinIO như ``S3 tại nhà'' và giải thích lỗi hai
endpoint. Chừng đó đủ để bài nộp hoạt động. Chương này hỏi thiết kế đó
tốn gì khi file trở nên nhiều, lớn và quan trọng --- khoảnh khắc bài
luận đã chấm của một học sinh là bản sao duy nhất ở bất kỳ đâu. Code
của dự án là mẫu vật: một \code{S3Config}, một \code{S3StorageService},
một cột tên \code{file\_object\_key}.

Ba câu hỏi sắp xếp chương này:

\begin{enumerate}
  \item \textbf{Byte nên sống ở đâu} --- trong cơ sở dữ liệu, trên đĩa,
  hay trong kho đối tượng --- và mỗi lựa chọn tốn chính xác những gì?
  \item \textbf{Ai nên chuyển byte} --- trình duyệt đi qua service, hay
  trình duyệt đi thẳng tới kho với một giấy phép có chữ ký?
  \item \textbf{Chuyện gì xảy ra khi hai kho bất đồng} --- một dòng gọi
  tên một object không tồn tại, hay một object không dòng nào gọi tên?
\end{enumerate}

\section{Vì sao kho đối tượng, và nó giải quyết gì}

Một bài nộp có hai nửa: các dữ kiện về nó (ai, khi nào, lần thứ mấy,
điểm bao nhiêu) và chính file. Dữ kiện có tính quan hệ và bé xíu; file
thì mờ đục và có thể nặng 20\,MB. File có ba nơi để ở.

\textbf{Trong Postgres, dưới dạng \code{bytea}.} Một kho, một
transaction, một bản backup. Cũng là nơi rắc rối bắt đầu. Mỗi lần đọc
dòng kéo cả blob qua shared buffers; \code{pg\_dump} lớn thêm bằng kích
thước mọi bài luận từng nộp; mỗi lần cập nhật ghi lại toàn bộ giá trị và
để lại một dead tuple cho autovacuum; write-ahead log chở mỗi byte hai
lần (WAL, rồi page); replication đẩy nó tới mọi standby; và một lượt tải
xuống phải hiện thân trọn vẹn trong JVM trước khi tới socket. Postgres
làm được, và tới vài gigabyte không ai để ý. Quá mức đó, mọi tác vụ vận
hành --- backup, restore, độ trễ replication, vacuum --- bị chi phối bởi
dữ liệu mà cơ sở dữ liệu chẳng thêm giá trị nào.

\textbf{Trên một volume mount vào service.} Nhanh và đơn giản cho tới
khi service có hai replica trên hai node, thứ đầu tiên Kubernetes muốn
làm. Một claim \code{ReadWriteOnce} (Chương~19) gắn với một node;
\code{ReadWriteMany} cần NFS hoặc một storage driver, và giờ service tự
sở hữu bố cục file, quyền, dọn dẹp và backup của riêng nó. Đĩa cũng nằm
trong vùng lỗi của pod.

\textbf{Trong kho đối tượng.} Không gian tên phẳng, API HTTP, độ bền
theo nội dung, và --- quyết định --- khả năng đưa cho client một URL rồi
bước ra khỏi đường dữ liệu. Backup thành một quy tắc replication của
bucket. Phục vụ file thành việc của CDN. Cơ sở dữ liệu chỉ giữ một key
500 ký tự.

\begin{decision}[một key trong Postgres, byte trong MinIO]
Dự án lưu metadata trong bảng \code{submission} và file dưới một key
trong bucket \code{ts-submissions}. Dòng tham chiếu object bằng tên;
không gì tham chiếu ngược từ object tới dòng. Các phương án khác thua ở
trục quan trọng với một trường học: bài nộp được ghi một lần, đọc vài
lần, và giữ nhiều năm --- hồ sơ mà kho đối tượng sinh ra để phục vụ và
cơ sở dữ liệu làm tệ nhất. Quyết định chỉ đảo ngược nếu file rất nhỏ
(vài KB), cần nguyên tử với dòng trong cùng transaction, và được đọc
mỗi lần truy cập dòng --- một ảnh đại diện thu nhỏ, chẳng hạn. Ngay cả
khi đó câu trả lời thật thà thường vẫn là ``vẫn S3, cộng một cache''.
\end{decision}

\section{Đọc code của dự án}

\subsection{Hai client, hai endpoint}

Cửa ngõ của mọi thứ là \code{S3Config}. Đọc nó với môi trường compose
trong đầu (\code{APP\_S3\_ENDPOINT=http://minio:9000},
\code{APP\_S3\_PUBLIC\_ENDPOINT=http://localhost:9000}):

\begin{lstlisting}[language=Java]
@Bean
public S3Client s3Client() {
    return S3Client.builder()
            .endpointOverride(URI.create(endpoint))            (*@\co{1}@*)
            .region(Region.of(region))
            .credentialsProvider(credentialsProvider())
            .serviceConfiguration(S3Configuration.builder()
                    .pathStyleAccessEnabled(true).build())     (*@\co{2}@*)
            .build();
}

@Bean
public S3Presigner s3Presigner() {
    return S3Presigner.builder()
            .endpointOverride(URI.create(publicEndpoint))      (*@\co{3}@*)
            .region(Region.of(region))
            .credentialsProvider(credentialsProvider())
            .serviceConfiguration(S3Configuration.builder()
                    .pathStyleAccessEnabled(true).build())
            .build();
}
\end{lstlisting}

\begin{description}
  \item[\co{1}] Client chuyển byte nói chuyện với tên \emph{nội bộ}.
  Trong Compose hay trong cụm, \code{minio} là một tên DNS trên mạng
  riêng (Chương~18); từ trình duyệt nó chẳng là gì.
  \item[\co{2}] Địa chỉ kiểu path-style: bucket nằm trong đường dẫn URL,
  nên request là \code{/ts-submissions/key} trên host \code{minio:9000}.
  Mặc định của AWS là kiểu virtual-host, bucket là một subdomain DNS ---
  host \code{ts-submissions.s3.amazonaws.com}, đường dẫn \code{/key}.
  MinIO cũng làm được, nhưng chỉ khi nó biết domain của chính mình
  (\code{MINIO\_DOMAIN}) và có wildcard DNS cho domain đó; trên mạng
  gia đình cả hai đều không có, nên dạng path là dạng duy nhất phân giải
  được.
  \item[\co{3}] Presigner không bao giờ gửi request. Nó chỉ tính chữ ký
  --- và phải tính cho host mà \emph{trình duyệt} sẽ liên hệ. Nó là một
  bean riêng chỉ vì lý do đó. Mục~\ref{sec:sigv4} giải thích vì sao
  không thể đổi host sau khi ký.
\end{description}

\subsection{Key, upload, download}

\code{S3StorageService} đủ ngắn để đọc trọn; các dòng quan trọng là
đây:

\begin{lstlisting}[language=Java]
@PostConstruct
void ensureBucketExists() {
    try {
        s3Client.headBucket(HeadBucketRequest.builder()
                .bucket(bucket).build());
    } catch (NoSuchBucketException e) {                     (*@\co{1}@*)
        s3Client.createBucket(CreateBucketRequest.builder()
                .bucket(bucket).build());
    }
}

public String upload(MultipartFile file, Long assignmentId,
                     Long studentId) {
    String key = "submissions/%d/%d/%s-%s".formatted(
            assignmentId, studentId, UUID.randomUUID(), original); (*@\co{2}@*)
    s3Client.putObject(PutObjectRequest.builder()
                    .bucket(bucket).key(key)
                    .contentType(file.getContentType())
                    .contentLength(file.getSize()).build(),
            RequestBody.fromInputStream(file.getInputStream(),
                    file.getSize()));                          (*@\co{3}@*)
    return key;
}

public String presignedDownloadUrl(String key) {
    GetObjectPresignRequest presignRequest = GetObjectPresignRequest
            .builder()
            .signatureDuration(Duration.ofMinutes(presignedExpiryMinutes))
            .getObjectRequest(GetObjectRequest.builder()
                    .bucket(bucket).key(key).build())
            .build();
    return s3Presigner.presignGetObject(presignRequest).url().toString();
}
\end{lstlisting}

\begin{description}
  \item[\co{1}] Chỉ \code{NoSuchBucketException} được bắt. Một
  connection refused --- MinIO chưa lên --- là
  \code{SdkClientException}, thoát ra ngoài, làm hỏng việc tạo bean, và
  kéo cả service xuống. Manifest Kubernetes của MinIO nói vậy ngay ở
  comment đầu: course-service sẽ \code{CrashLoopBackOff} cho tới khi
  MinIO trả lời. Compose che điều này bằng
  \code{depends\_on: service\_healthy}.
  \item[\co{2}] Key là toàn bộ mô hình dữ liệu của một kho đối tượng.
  S3 không có thư mục; dấu gạch chéo chỉ là ký tự trong chuỗi. Nhưng
  \emph{tiền tố} là thứ bạn liệt kê theo, nên
  \code{submissions/\{assignmentId\}/\{studentId\}/} nghĩa là ``mọi lần
  nộp của học sinh này cho bài này'' là một lời gọi \code{ListObjectsV2}
  với một tiền tố. UUID bảo đảm key mới cho mỗi lần nộp, nên không gì
  bị ghi đè và versioning không cần thiết cho an toàn. Tên file gốc
  được nối vào để dễ đọc trong console; cơ sở dữ liệu giữ riêng nó trong
  \code{file\_name} vì key không phải chỗ tốt cho dấu cách, Unicode, hay
  \code{..}.
  \item[\co{3}] Byte chảy \emph{xuyên qua} JVM này: hỗ trợ multipart
  của Spring đã parse chúng (đổ ra file tạm khi vượt một ngưỡng nhỏ),
  và giờ chúng được stream tới MinIO. Thread request bị chiếm suốt
  lượt truyền, hai lần.
\end{description}

Và đường nối với nửa quan hệ, trong \code{SubmissionService}:

\begin{lstlisting}[language=Java]
@Transactional
public SubmissionResponse submit(...) {
    ...
    if (hasFile) {
        String key = s3StorageService.upload(file, assignmentId,
                student.getId());                              (*@\co{1}@*)
        submission.setFileObjectKey(key);
        submission.setFileName(file.getOriginalFilename());
    }
    submission = submissionRepository.save(submission);        (*@\co{2}@*)
    return toResponse(submission);
}
\end{lstlisting}

\begin{description}
  \item[\co{1}] Một lời gọi HTTP tới hệ thống khác, bên trong một
  transaction cơ sở dữ liệu. Nó không, và không thể, là một phần của
  transaction đó.
  \item[\co{2}] Nếu insert này thất bại --- ràng buộc unique trên
  \code{(assignment\_id, student\_id, attempt\_number)} thua một cuộc
  đua giữa hai tab, hay Postgres thoáng không với tới được ---
  transaction rollback và object đã upload ở \co{1} nằm lại trong
  bucket, mang một key không dòng nào sẽ ghi lại. Chương~11 gọi đây là
  bài toán dual-write khi hệ thống thứ hai là Kafka; vẫn bài toán đó với
  một hệ thống thứ hai khác.
\end{description}

\section{Signature Version 4, và vì sao host được hàn chết vào}
\label{sec:sigv4}

Một presigned URL là một request S3 bình thường với phần xác thực
chuyển từ header sang chuỗi truy vấn. Link tải xuống của dự án trông
thế này (đã ngắt dòng):

\begin{lstlisting}
http://localhost:9000/ts-submissions/submissions/7/3/2f9c...-essay.pdf
  ?X-Amz-Algorithm=AWS4-HMAC-SHA256
  &X-Amz-Credential=minioadmin/20260907/us-east-1/s3/aws4_request
  &X-Amz-Date=20260907T091500Z
  &X-Amz-Expires=900
  &X-Amz-SignedHeaders=host
  &X-Amz-Signature=3c1f...e9a0
\end{lstlisting}

Chữ ký được tính qua hai giai đoạn. Trước hết một \emph{khóa ký} được
dẫn xuất từ secret qua một chuỗi HMAC trên ngày, region, tên dịch vụ và
chuỗi cố định \code{aws4\_request}, nên bản thân secret không bao giờ
chạm vào request. Rồi khóa đó ký một \emph{string to sign} chứa SHA-256
của một \emph{canonical request}: phương thức HTTP, URI chuẩn hóa
(bucket và key, kiểu path), chuỗi truy vấn đã sắp xếp, các header được
ký, và mã băm của payload. Nhìn \code{X-Amz-SignedHeaders=host}: header
\code{Host} luôn được ký. Đổi host --- \code{minio} thành
\code{localhost}, hay \code{localhost} thành \code{127.0.0.1} --- là
canonical request đổi, mã băm đổi, chữ ký không khớp nữa, và MinIO trả
lời \code{403 SignatureDoesNotMatch}. Không có cách nào ký với tên nội
bộ rồi viết lại link cho thế giới bên ngoài. Ai ký phải biết chính xác
origin mà trình duyệt sẽ dùng.

\begin{handson}[cố tình làm hỏng một chữ ký]
Đăng nhập, nộp một bài có file, và chép \code{downloadUrl} mà API trả
về. Nó hoạt động đúng như được phát:
\begin{lstlisting}[language=bash]
$ curl -s -o /dev/null -w "%{http_code}\n" "$URL"
200
\end{lstlisting}
Giờ chỉ đổi host, sang một địa chỉ tới đúng container MinIO đó:
\begin{lstlisting}[language=bash]
$ curl -s "${URL/localhost/127.0.0.1}" | head -c 160
<?xml version="1.0" encoding="UTF-8"?><Error>
<Code>SignatureDoesNotMatch</Code><Message>The request signature we
calculated does not match the signature you provided.
\end{lstlisting}
Cùng server, cùng object, cùng thông tin xác thực --- và một header
\code{Host} khác là đủ. Đó là toàn bộ lý do \code{public-endpoint} tồn
tại.
\end{handson}

\begin{pitfall}[link tải xuống chạy trong Compose và hỏng trong cụm]
Triệu chứng: trên cụm, bấm vào file của một bài nộp mở ra
\code{http://localhost:9000/...} và trình duyệt báo connection refused.
Nguyên nhân: Deployment của course-service đặt \code{APP\_S3\_ENDPOINT}
là \code{http://minio:9000} nhưng không bao giờ đặt
\code{APP\_S3\_PUBLIC\_ENDPOINT}, nên giá trị mặc định
\code{http://localhost:9000} của config-server được ký vào mọi link ---
một địa chỉ từng nghĩa là ``laptop của lập trình viên'' và giờ nghĩa là
``laptop của người dùng''. Sửa hai phần: mở cổng S3 của MinIO qua một
quy tắc Ingress trên host riêng của nó (chữ ký path-style bao gồm cả
đường dẫn, nên Ingress không được thêm tiền tố), và đặt public endpoint
thành origin đó:
\begin{lstlisting}[language=yaml]
- { name: APP_S3_PUBLIC_ENDPOINT,
    value: "https://files.ts-server.tailbfe002.ts.net" }
\end{lstlisting}
Chừng nào chưa có hostname thứ hai, giải pháp tạm thật thà là stream
lượt tải qua course-service (\code{GetObject} vào body của response) để
trình duyệt không bao giờ cần với tới MinIO --- chậm hơn, nhưng đúng.
\end{pitfall}

\section{Ai chuyển byte}

Đường upload hôm nay là: trình duyệt $\rightarrow$ Traefik $\rightarrow$
api-gateway $\rightarrow$ course-service $\rightarrow$ MinIO. Mỗi chặng
thấy mỗi byte. Service cưỡng chế \code{max-file-size: 20MB} và
\code{max-request-size: 25MB} trong cấu hình multipart, nên trường hợp
xấu nhất có giới hạn; nhưng mỗi lượt upload giữ một thread Tomcat và một
file tạm suốt thời gian của nó, gateway phản ứng (Chương~4) chuyển tiếp
body từng khối nhưng vẫn tiêu event loop cho việc đó, và mọi giới hạn
body hay idle timeout ở tầng ingress áp lên trước tất cả. Dưới hạn nộp
--- 200 học sinh upload lúc 23:59 --- thread pool của service là nút
thắt cho công việc nó không cần làm.

Phương án thay thế là mẫu Chương~9 nhắc tới cho tải xuống, áp dụng cho
tải lên: service ký một \emph{PUT} cho một key nó chọn, trình duyệt
upload thẳng tới MinIO, và chỉ khi đó service mới ghi dòng. Service tốn
vài micro giây mỗi file và không bao giờ thấy byte. Cái nó từ bỏ: không
kiểm tra được nội dung trước khi nó hạ cánh (quét virus, nhận diện
loại) và phải biết, từ đâu đó, rằng upload thực sự hoàn tất.
Mục~\ref{sec:presigned-put} xây dựng điều này cho tử tế.

Quá khoảng 100\,MB, một PUT đơn là công cụ sai trên mọi mạng:
\emph{multipart upload} chia object thành các phần ít nhất 5\,MB,
upload song song, tiếp tục sau sự cố bằng cách gửi lại một phần, và ráp
chúng phía server trong một lời gọi cuối. \code{S3TransferManager} của
SDK tự làm việc này trên một ngưỡng. Cái đuôi vận hành: upload bị bỏ dở
để lại các phần vô hình vẫn bị tính tiền và đếm --- mọi bucket nên mang
một quy tắc lifecycle hủy multipart upload dang dở sau vài ngày.

\section{Ưu và nhược, thẳng thắn}

\begin{center}
\begin{tabular}{@{}p{0.3\linewidth}p{0.62\linewidth}@{}}
\toprule
\textbf{Điểm mạnh} & \textbf{Dự án trả giá ở đâu} \\
\midrule
Cơ sở dữ liệu giữ nhỏ; backup chỉ là metadata &
  Hai kho phải backup, và chỉ một trong số đó có transaction \\
API chuẩn: MinIO hôm nay, S3 ngày mai với hai biến môi trường &
  Hai endpoint phải giữ cho đúng, và cụm đặt sai một \\
Presigned download đưa service ra khỏi đường dữ liệu &
  Upload vẫn đi qua service, mỗi file một thread \\
Key cho ``thư mục'' miễn phí theo tiền tố &
  Thiết kế key là cánh cửa một chiều; đổi tên nghĩa là copy \\
Độ bền là việc của kho &
  Trên phần cứng này nó là một container và một claim 5\,GB \\
Thông tin xác thực là hai chuỗi &
  Đó là thông tin \emph{root} của MinIO, trong mọi pod course-service \\
\bottomrule
\end{tabular}
\end{center}

Dòng cuối đáng nhấn mạnh. \code{minioadmin} tạo được bucket, xóa được
bucket, và đọc được mọi object trong mọi bucket. Service cần
\code{PutObject}, \code{GetObject} và \code{HeadBucket} trên một bucket.
Quyền tối thiểu ở đây là một user MinIO với policy giới hạn vào
\code{arn:aws:s3:::ts-submissions/*}; trên AWS là một IAM role gắn với
ServiceAccount của pod (IRSA) để không tồn tại khóa tĩnh nào. Người
phỏng vấn hỏi ``service của bạn xác thực với S3 thế nào?'' đang nghe
xem bạn nói ``một IAM role'' hay ``một key trong biến môi trường'' ---
và, nếu là vế sau, bạn có biết đó là câu trả lời yếu hơn không.

\section{Chuyện gì gãy}

\begin{pitfall}[service không chịu khởi động khi thiếu bucket]
Triệu chứng: course-service \code{CrashLoopBackOff} trên cụm mới; log
kết thúc bằng \code{Error creating bean with name 's3StorageService'}
với một \code{Connection refused} nằm sâu vài tầng nguyên nhân. Nguyên
nhân: kiểm tra bucket trong \code{@PostConstruct} coi mọi thất bại
ngoài ``thiếu bucket'' là chí mạng, mà pod MinIO vẫn đang kéo image.
Back-off khởi động lại của kubelet cuối cùng rơi vào lúc MinIO sẵn sàng,
nên nó tự lành --- chậm. Sửa: bắt cả \code{SdkClientException}, log
cảnh báo, và để readiness probe (Chương~17) giữ lưu lượng tránh xa cho
tới lần \code{headBucket} thành công đầu tiên; một phụ thuộc lưu trữ
nên làm suy giảm tính năng file, không phải cả service.
\end{pitfall}

\begin{pitfall}[object mồ côi]
Triệu chứng: bucket chứa nhiều object hơn số key khác null trong bảng
\code{submission}; không ai để ý suốt một năm; claim 5\,GB đầy. Nguyên
nhân: dual write trong \code{submit}. Upload trước, insert sau: insert
hỏng là object mồ côi. Insert trước, upload sau: upload hỏng để lại
dòng trỏ vào khoảng không, và học sinh thấy ``đã đính kèm file'' với
một link chết. Thứ tự nào cũng không tự an toàn. Cách giải chuẩn là làm
một phía lũy đẳng và phía kia đối chiếu được: insert dòng với trạng
thái \code{PENDING\_FILE} và key đã chọn, upload, rồi lật trạng thái
sang \code{SUBMITTED} trong transaction thứ hai; một job hằng đêm xóa
object dưới các key mà dòng vẫn \code{PENDING\_FILE} sau một giờ, và
liệt kê tiền tố để tìm key không dòng nào gọi tên.
Mục~\ref{sec:presigned-put} gấp điều này vào luồng mới, nơi nó gần như
miễn phí.
\end{pitfall}

\section{Chuyện gì gãy lúc 3 giờ sáng}

Pod MinIO biến mất --- node khởi động lại, claim chưa gắn lại. Đi qua
từng endpoint:

\begin{itemize}
  \item \textbf{Bài nộp dạng text và link} vẫn chạy: \code{upload} chỉ
  được gọi khi có file, và \code{save} chỉ chạm Postgres.
  \item \textbf{Bài nộp có file} thất bại với 500 sau ngân sách retry
  của SDK --- vài giây giữ thread mỗi lượt. Dưới tải, các thread này
  tích tụ và bỏ đói cả các bài nộp text lẽ ra đã thành công. Đây là lập
  luận của Chương~17 cho một bulkhead: một executor riêng, có giới hạn,
  cho các lời gọi lưu trữ.
  \item \textbf{Liệt kê bài nộp} chạy và thậm chí trả về link tải
  xuống. Ký là thuần tính toán; không request nào được gửi. Link chỉ
  hỏng khi bấm, và đó là cách hỏng đúng --- metadata là thật, byte tạm
  thời không sẵn có.
  \item \textbf{Khởi động lại course-service} trong lúc sự cố thì không
  lên được, vì lý do trong cạm bẫy ở trên. Một sự cố của kho file đã
  thành sự cố của chấm điểm, ghi danh, và mọi thứ khác service làm.
\end{itemize}

Giờ chiều ngược lại: Postgres chết và MinIO bình thường. Mọi thứ thất
bại, kể cả các upload lẽ ra đã thành công ở tầng lưu trữ --- và mỗi
upload chạm tới \co{1} trong \code{submit} rồi hỏng ở \co{2} để lại một
object mồ côi. Một cú chớp cơ sở dữ liệu lúc 3 giờ sáng sinh ra một
đống nhỏ object vô chủ, và đó là lý do cụ thể job đối chiếu tồn tại.

\section{Nâng cấp giải pháp: presigned PUT}
\label{sec:presigned-put}

Bản nâng cấp đưa byte ra khỏi service và đóng cửa sổ mồ côi cùng lúc.
Ba lời gọi thay cho một:

\begin{enumerate}
  \item \code{POST /assignments/\{id\}/submissions/upload-url} ---
  service kiểm tra học sinh được nộp, chọn key, ký một PUT, và insert
  một dòng \code{PENDING\_FILE}.
  \item Trình duyệt PUT file lên MinIO với URL đó.
  \item \code{POST /assignments/\{id\}/submissions} kèm key --- service
  xác nhận object tồn tại bằng một \code{HEAD}, rồi hoàn tất dòng.
\end{enumerate}

Storage service có thêm hai phương thức. Cả hai đều trọn vẹn:

\begin{lstlisting}[language=Java]
/** Sign a PUT for a key we chose; the browser must send exactly this
 *  content type, because it is part of the signature. */
public String presignedUploadUrl(String key, String contentType,
                                 long contentLength) {
    PutObjectRequest put = PutObjectRequest.builder()
            .bucket(bucket)
            .key(key)
            .contentType(contentType)
            .contentLength(contentLength)                  (*@\co{1}@*)
            .build();
    PutObjectPresignRequest req = PutObjectPresignRequest.builder()
            .signatureDuration(Duration.ofMinutes(10))
            .putObjectRequest(put)
            .build();
    return s3Presigner.presignPutObject(req).url().toString();
}

/** Confirm an object exists and report its size; empty if missing. */
public Optional<Long> objectSize(String key) {
    try {
        HeadObjectResponse head = s3Client.headObject(
                HeadObjectRequest.builder().bucket(bucket).key(key).build());
        return Optional.of(head.contentLength());
    } catch (NoSuchKeyException e) {                        (*@\co{2}@*)
        return Optional.empty();
    }
}
\end{lstlisting}

\begin{description}
  \item[\co{1}] Ký \code{Content-Length} và \code{Content-Type} biến
  chúng thành hợp đồng: trình duyệt gửi kích thước hay loại khác sẽ nhận
  \code{403}. Đó là cách giới hạn 20\,MB sống sót khi rời khỏi service.
  \item[\co{2}] \code{HEAD} tốn một vòng đi-về và không byte nào. Đó là
  bằng chứng rẻ nhất có thể rằng bước~2 đã xảy ra.
\end{description}

Submission service tách phương thức cũ làm hai. Tiền tố key được dựng
từ dữ kiện phía server, nên client không thể nhận vơ upload của người
khác:

\begin{lstlisting}[language=Java]
public UploadTicket prepareUpload(Long assignmentId, Long studentUserId,
                                  String fileName, String contentType,
                                  long size) {
    Assignment assignment = requireOpenAssignment(assignmentId);
    Student student = requireEnrolled(assignment, studentUserId);
    if (size > MAX_FILE_BYTES) {
        throw new ApiException(413, "File exceeds 20 MB");
    }
    String key = "submissions/%d/%d/%s-%s".formatted(
            assignmentId, student.getId(), UUID.randomUUID(),
            safeName(fileName));                            (*@\co{1}@*)
    Submission pending = Submission.builder()
            .assignment(assignment).student(student)
            .attemptNumber(nextAttempt(assignmentId, student.getId()))
            .fileObjectKey(key).fileName(fileName)
            .status("PENDING_FILE")                         (*@\co{2}@*)
            .build();
    submissionRepository.save(pending);
    String url = s3StorageService.presignedUploadUrl(key, contentType,
            size);
    return new UploadTicket(pending.getId(), key, url);
}

public SubmissionResponse completeUpload(Long submissionId,
                                         Long studentUserId) {
    Submission s = findOrThrow(submissionId);
    if (!s.getStudent().getUserId().equals(studentUserId)
            || !"PENDING_FILE".equals(s.getStatus())) {
        throw new ApiException(409, "Nothing to complete");
    }
    long size = s3StorageService.objectSize(s.getFileObjectKey())
            .orElseThrow(() -> new ApiException(400,
                    "Upload not found; request a new upload URL")); (*@\co{3}@*)
    if (size > MAX_FILE_BYTES) {
        throw new ApiException(413, "File exceeds 20 MB");
    }
    s.setStatus("SUBMITTED");
    return toResponse(submissionRepository.save(s));
}
\end{lstlisting}

\begin{description}
  \item[\co{1}] \code{safeName} bỏ dấu phân cách đường dẫn và ký tự
  ngoài ASCII; tên hiển thị ở lại trong \code{file\_name}.
  \item[\co{2}] Dòng tồn tại \emph{trước} byte. Crash ở bất kỳ đâu sau
  điểm này để lại một dòng \code{PENDING\_FILE} mà job dọn dẹp tìm
  được, thay vì một object không ai tìm được.
  \item[\co{3}] Con đường duy nhất tới \code{SUBMITTED} là qua một
  \code{HEAD} thành công. Dòng không bao giờ nhận có một file không ở
  đó.
\end{description}

Job dọn dẹp giờ là một truy vấn, không phải một lượt quét bucket:

\begin{lstlisting}[language=Java]
@Scheduled(cron = "0 15 * * * *")   // hourly, at :15
public void reapAbandonedUploads() {
    List<Submission> stale = submissionRepository
            .findByStatusAndSubmittedAtBefore("PENDING_FILE",
                    LocalDateTime.now().minusHours(1));
    for (Submission s : stale) {
        s3StorageService.deleteQuietly(s.getFileObjectKey()); // idempotent
        submissionRepository.delete(s);
    }
}
\end{lstlisting}

MinIO phải cho phép PUT cross-origin từ trình duyệt. Quy tắc CORS của
bucket được đặt một lần bằng MinIO client:

\begin{lstlisting}[language=bash]
$ mc cors set local/ts-submissions cors.json
# cors.json: allow PUT from the frontend origin, expose ETag
{ "CORSRules": [ { "AllowedOrigins": ["http://localhost:3000"],
    "AllowedMethods": ["PUT", "GET"],
    "AllowedHeaders": ["Content-Type", "Content-Length"],
    "ExposeHeaders": ["ETag"], "MaxAgeSeconds": 3600 } ] }
\end{lstlisting}

Và phía frontend là một \code{fetch} trơn; không instance axios, không
gateway, không JWT --- URL \emph{chính là} thông tin xác thực:

\begin{lstlisting}
const ticket = await api.post(
  `/assignments/${id}/submissions/upload-url`,
  { fileName: file.name, contentType: file.type, size: file.size });

const put = await fetch(ticket.url, {
  method: 'PUT',
  headers: { 'Content-Type': file.type },   // must match the signature
  body: file,
});
if (!put.ok) throw new Error(`upload failed: ${put.status}`);

await api.post(`/submissions/${ticket.submissionId}/complete`);
\end{lstlisting}

\begin{decision}[presigned PUT so với upload qua proxy]
Đi qua proxy khi service phải thấy byte trước khi lưu --- quét virus,
mã hóa lại ảnh, nhận diện content-type --- hoặc khi client không thể
với tới endpoint lưu trữ. Presign khi file lớn hoặc nhiều và service
chẳng thêm gì cho lượt truyền: thread pool, giới hạn body ở ingress, và
chặng mạng kép đều biến mất. Trường hợp của dự án là vế thứ hai, với
một lưu ý quyết định thứ tự công việc: presigned upload đòi một origin
MinIO trình duyệt với tới được, thứ cụm chưa có. Cạm bẫy ở trên và bản
nâng cấp này chung một điều kiện tiên quyết --- một hostname công khai
cho file --- và nên được lên lịch cùng nhau.
\end{decision}

\section{Ngoài dự án}

Trên AWS, cùng đoạn code chạy với S3 khi bỏ endpoint override và thay
khóa tĩnh bằng IAM role; CloudFront signed URL đóng vai presigner với
một CDN đứng trước, và \emph{S3 POST policy} là người anh em cũ hơn,
thân thiện với form trình duyệt, của presigned PUT, có sẵn điều kiện
\code{content-length-range}. Google Cloud Storage và Azure Blob Storage
có cùng ba ý tưởng dưới tên khác --- signed URL, SAS token, resumable
upload. MinIO production chạy dạng phân tán với erasure coding qua nhiều
ổ đĩa và node; một container với một claim là hình hài để phát triển,
và SDK không phân biệt được, đó chính là ý nghĩa của chuẩn. Phiên bản
phỏng vấn của chương này là một câu hỏi --- ``người dùng upload một
video 2\,GB; dẫn tôi đi qua'' --- và câu trả lời senior nêu presigned
hoặc multipart upload, hợp đồng kích thước và loại trong chữ ký, một
bước xác nhận, một job dọn mồ côi, và một quy tắc lifecycle cho các
phần dang dở, theo gần đúng thứ tự đó.

\section{Câu hỏi từ phỏng vấn thật}

\subsection*{Q: Người dùng bắt đầu upload video 1\,GB và service hết bộ nhớ. Sửa đi.}

Trước hết xác định byte là vấn đề, không phải metadata: một heap dump
giữa lúc upload cho thấy các buffer multipart hay một
\code{ByteArrayOutputStream}, và giới hạn bộ nhớ của pod (Chương~17) là
thứ giết nó. Sửa ngắn hạn là bảo đảm Spring đổ ra đĩa
(\code{file-size-threshold} nhỏ, một volume tạm đủ chỗ) và stream tới
kho với \code{Content-Length} biết trước, như
\code{S3StorageService.upload} đã làm --- đưa bạn tới hàng trăm MB mà
heap không phình, nhưng vẫn một thread và một file tạm mỗi upload. Sửa
thật là đưa service ra khỏi đường đi: presigned PUT cho mọi thứ dưới
khoảng 100\,MB, và multipart upload với \code{S3TransferManager} hoặc
chia phần phía trình duyệt cho phần trên, kèm quy tắc lifecycle hủy
upload dang dở sau bảy ngày. Khi đó bộ nhớ của service độc lập với kích
thước file theo cấu trúc, và \code{max-file-size} trở thành một
\code{Content-Length} đã ký thay vì một cấu hình servlet.

\subsection*{Q: Chúng ta cần phục vụ mười nghìn lượt tải xuống mỗi giây. Làm sao, mà không chạm vào service?}

Presigned GET đã nghĩa là service không làm I/O nào mỗi lượt tải, chỉ
một chữ ký --- vài micro giây, không mạng. Mười nghìn mỗi giây vì thế
là câu hỏi về kho và biên. Kho đối tượng mở rộng đọc theo chiều ngang
bằng cách phân bố tiền tố key; một CDN phía trước (CloudFront với signed
URL, hay CDN nào kéo từ origin là bucket) phục vụ object lặp lại từ
cache và hấp thụ fan-out. Giữ việc ký rẻ bằng cách không cache gì ---
chữ ký là phi trạng thái --- và đặt hạn dùng khớp cửa sổ cache của CDN
để một object đã cache không sống lâu hơn quyền của nó bao nhiêu. Cái
bẫy là ký trong mỗi lời gọi liệt kê: \code{toResponse} trong dự án ký
cho mỗi dòng của mỗi danh sách, ổn với hạn 15 phút nhưng nghĩa là một
danh sách nghìn bài nộp tính nghìn chuỗi HMAC. Ký lười, ở một endpoint
``lấy link'' riêng, khi lưu lượng biện minh cho nó.

\subsection*{Q: Object có trong bucket nhưng dòng thì thiếu, hoặc dòng có mà object không. Đối chiếu thế nào?}

Nói trước rằng cả hai là kết quả tất yếu của dual write và cách sửa là
thiết kế, không phải một script. Thiết kế: ghi dòng ở trạng thái
\code{PENDING} với key trước khi upload, xác nhận bằng \code{HEAD} trước
khi đánh dấu hoàn tất, và không bao giờ để dòng nhận một file nó chưa
xác minh. Đối chiếu khi đó có hai nửa rẻ. Dòng kẹt ở \code{PENDING} lâu
hơn hạn presign bị xóa cùng object của chúng --- một truy vấn, không
phải quét. Object không dòng được tìm bằng cách liệt kê tiền tố và
anti-join với bảng, việc bạn kham nổi hằng đêm vì tiền tố là theo bài
tập; trên AWS một báo cáo S3 Inventory thay cho lượt liệt kê. Làm gì
thì làm, hãy để lệnh xóa object lũy đẳng để job chạy lại được.

\subsection*{Q: Bảo vệ một presigned URL thế nào?}

Bốn núm vặn, và câu trả lời senior gọi tên cả bốn. \emph{Phạm vi}: ký
một phương thức trên một key; URL GET không dùng để PUT được, và URL
PUT không dùng cho key khác được vì URI chuẩn hóa được ký. \emph{Hạn
dùng}: phút, không phải ngày; dự án dùng 15 cho tải xuống và 10 là đủ
cho tải lên --- đủ dài cho kết nối chậm, đủ ngắn để một link rò rỉ
trong log chat đã chết khi ai đó để ý. \emph{Hợp đồng nội dung}: ký
\code{Content-Type} và \code{Content-Length} khi upload để client không
đổi được điều nó đã khai. \emph{Danh tính người ký}: thông tin xác thực
đằng sau chữ ký nên là một user quyền tối thiểu trên đúng bucket đó, để
ngay cả người ký bị chiếm cũng không liệt kê hay xóa được. Và nói điều
người ta hay quên: presigned URL là một bearer credential --- không log
ở đâu, chỉ trả về qua HTTPS, và không bao giờ đặt vào một redirect mà
proxy có thể cache.

\subsection*{Q: Chuyển từ MinIO sang AWS S3 không gián đoạn.}

Vì code nói API S3, cuộc di cư là một lần chuyển dữ liệu cộng một thay
đổi cấu hình, và mẹo nằm ở thứ tự. Thiết lập bucket replication hoặc
chạy \code{mc mirror --watch} từ MinIO sang bucket S3 cho tới khi độ
trễ chỉ còn vài giây. Chuyển \emph{ghi} trước bằng cách đổi
\code{APP\_S3\_ENDPOINT} và thông tin xác thực (giờ là một IAM role)
trên một rolling restart --- object mới hạ cánh ở S3, mirroring tiếp
tục chép bất kỳ thứ gì còn hạ cánh ở MinIO từ các pod chưa khởi động
lại. Chuyển \emph{public endpoint} sau cùng, sau một lượt mirror cuối,
để link tải xuống chỉ bao giờ được ký cho host có object. Key không
đổi, nên cơ sở dữ liệu không cần migration. Giữ MinIO chỉ-đọc một tuần
làm phương án lùi, rồi ngừng. Thứ gãy là cách địa chỉ: kiểu virtual-host
trở nên khả dụng, kiểu path vẫn chạy, và đoạn code nào ghép URL bằng tay
thay vì hỏi SDK chính là nơi bạn sẽ tìm thấy lỗi.

\subsection*{Q: Chi phí lưu trữ tăng mỗi tháng. Làm gì?}

Đo theo tiền tố trước --- bài tập nào, năm nào. Rồi dùng công cụ của
chính kho thay vì logic ứng dụng: một quy tắc lifecycle chuyển object
cũ hơn một học kỳ sang tầng ít truy cập, và cũ hơn chính sách lưu giữ
sang lưu trữ lạnh hoặc xóa; một quy tắc hủy multipart upload dang dở,
thứ vô hình mà vẫn tính tiền; và tắt versioning, hoặc có hạn, trừ khi
nhu cầu tuân thủ nói khác. Trong ứng dụng, ngừng lưu thứ tái tạo được
(thumbnail, bản chuyển mã) và khử trùng lặp theo mã băm nội dung nếu
học sinh upload cùng một mẫu. Ở dự án này con số nhỏ --- 250 lượt
upload 20\,MB là đầy claim 5\,GB --- nhưng cùng bộ quy tắc đó là thứ
giữ hóa đơn S3 phẳng ở quy mô gấp nghìn lần.

\begin{quiz}
\begin{enumerate}
  \item Hai bean, \code{S3Client} và \code{S3Presigner}, chung thông tin
  xác thực và region, khác nhau một dòng. Nêu dòng đó, và giải thích
  chuyện gì sai nếu presigner dùng endpoint nội bộ.
  \item Một đồng nghiệp đề xuất giữ một endpoint và để gateway viết lại
  \code{minio:9000} thành host công khai trong JSON trả về. Dùng cấu
  trúc chữ ký Version 4, giải thích vì sao link đã viết lại trả về 403.
  \item Trong \code{SubmissionService.submit}, upload diễn ra bên trong
  một phương thức \code{@Transactional}. Nêu transaction bao phủ và
  không bao phủ điều gì, và mô tả thứ để lại khi \code{save} thất bại.
  \item Bố cục key là
  \code{submissions/\{assignment\}/\{student\}/\{uuid\}-\{name\}}. Nêu
  một thao tác bố cục này làm rẻ và một thao tác nó làm đắt, và nói vì
  sao UUID khiến versioning của bucket không cần thiết cho an toàn.
  \item MinIO chết lúc 3 giờ sáng. Trong các việc này, việc nào vẫn
  thành công: nộp câu trả lời text, liệt kê bài nộp, mở link tải xuống,
  khởi động lại course-service? Giải thích mỗi việc trong một mệnh đề.
  \item Trong luồng presigned PUT, kẻ tấn công lấy được một URL upload
  hợp lệ từ phiên của chính mình và cố gắn object đó vào bài nộp của
  học sinh khác. Chỉ ra hai kiểm tra chặn họ.
\end{enumerate}
\end{quiz}
